% القسم: المناهج
% الفصل: البراهين
% المبحث: البرهان بالتناقض

\documentclass[../../../include/open-logic-section]{subfiles}

\begin{document}

\olfileid{mth}{prf}{con}

\olsection{البرهان بالتناقض}

البرهان بالتناقض، في المقام الأول، نمط استدلال يستعمل لإثبات القضايا
السالبة. افترض أنك تريد بيان أن قضية ما~$p$ \emph{كاذبة}، أي إنك تريد
بيان~$\lnot p$. والاستراتيجية الأرجح نجاحاً هي: (أ) أن تفترض صدق $p$،
و(ب) أن تبين أن هذا الافتراض يؤدي إلى شيء تعرف كذبه. وقد يكون «الشيء
المعروف كذبه» نتيجة تتعارض---تتناقض---مع $p$ نفسها، أو مع فرض آخر من
فروض القضية الكلية التي تبحثها. فمثلاً، يقتضي برهان «إذا كان $q$ فإن
$\lnot p$» افتراض صدق $q$ وإثبات~$\lnot p$ منه. وإذا أثبتت
$\lnot p$ بالتناقض، فهذا يعني افتراض $p$ إلى جانب~$q$. فإذا استطعت
إثبات $\lnot q$ من $p$، فقد بينت أن الفرض~$p$ يؤدي إلى شيء يناقض فرضك
الآخر~$q$، إذ لا يمكن أن يصدق $q$ و$\lnot q$ معاً. وبالطبع، لا بد أن
تستعمل في إثبات التناقض أنماط استدلال أخرى وأن تفك التعريفات. فلننظر في
مثال.

\begin{prop}
  إذا كان $A \subseteq B$ و$B = \emptyset$، فلا توجد في $A$
  !!{element}.
\end{prop}

\begin{proof}
  افترض أن $A \subseteq B$ و$B = \emptyset$. نريد أن نبين أنه لا توجد
  في $A$ !!{element}.
  \begin{quote}
    بما أن هذه قضية شرطية، نفترض مقدمها ونريد إثبات تاليها. والتالي هو:
    لا توجد في $A$ !!{element}. ويمكننا التصريح بذلك أكثر: ليس صحيحاً
    أنه يوجد~$x \in A$.
  \end{quote}
  لا توجد في $A$ !!{element} إذا وفقط إذا لم يكن صحيحاً أنه يوجد~$x$
  بحيث $x \in A$.
  \begin{quote}
    لقد حددنا إذن أن ما نريد إثباته قضية سالبة~$\lnot p$، وهي: ليس
    صحيحاً أنه يوجد $x \in A$. ولاستعمال البرهان بالتناقض علينا افتراض
    القضية الموجبة المقابلة~$p$، أي إنه يوجد $x \in A$، ثم اشتقاق
    تناقض منها. ونشير إلى أننا نقيم برهاناً بالتناقض بقولنا «للوصول إلى
    تناقض، افترض» أو حتى «افترض العكس» فحسب، ثم نذكر الفرض~$p$.
  \end{quote}
  افترض العكس: يوجد $x \in A$.
  \begin{quote}
    هذا هو الفرض الجديد الذي سنستعمله للحصول على تناقض. ولدينا فرضان
    آخران: $A \subseteq B$ و$B = \emptyset$. يعطينا الأول
    $x \in B$:
  \end{quote}
  بما أن $A \subseteq B$، فإن $x \in B$.
  \begin{quote}
    لكن بما أن $B = \emptyset$، يجب أن تكون كل !!{element} في $B$،
    ومنها $x$، أيضاً !!a{element} في~$\emptyset$.
  \end{quote}
  بما أن $B = \emptyset$، فإن $x \in \emptyset$. وهذا تناقض، لأن
  $\emptyset$ لا تحتوي، بحسب التعريف، على أي !!{element}.
  \begin{quote}
    بهذا يكتمل البرهان بالفعل: وصلنا من الافتراضات التي وضعناها إلى ما
    نحتاج إليه، وهو تناقض، وهذا يعني أنه لا يمكن أن تصدق الافتراضات كلها.
    ولما كان الافتراضان الأولان، $A \subseteq B$ و$B = \emptyset$، غير
    متنازع فيهما، فلا بد أن يكون آخر افتراض أدخلناه، وهو وجود
    $x \in A$، كاذباً. لكن يمكننا التصريح بذلك إن أردنا الدقة.
  \end{quote}
  إذن يجب أن يكون افتراضنا وجود $x \in A$ كاذباً، ومن ثم لا توجد في
  $A$ !!{element}، بالبرهان بالتناقض.
\end{proof}

كل قضية موجبة تكافئ، على نحو تافه، قضية سالبة: $p$ إذا وفقط إذا كان
$\lnot\lnot p$. ولذلك يمكن أيضاً استعمال البرهان بالتناقض لإثبات
القضايا الموجبة «بطريق غير مباشر» كما يأتي: لإثبات~$p$، اقرأها بوصفها
القضية السالبة $\lnot\lnot p$. فإذا استطعنا اشتقاق تناقض من
$\lnot p$، فقد أثبتنا $\lnot\lnot p$ بالتناقض، ومن ثم أثبتنا~$p$.

كنا نهدف في المثال الأخير إلى إثبات قضية سالبة، هي أنه لا توجد في $A$
!!{element}، ولذلك كان الافتراض الذي وضعناه بغرض البرهان بالتناقض، أي
وجود $x \in A$، قضية موجبة. فقد منحنا شيئاً نعمل عليه، هو العنصر
الافتراضي $x \in A$، الذي واصلنا الاستدلال بشأنه حتى وصلنا إلى
$x \in \emptyset$.

عندما تثبت قضية موجبة بطريق غير مباشر، يكون الافتراض الذي تضعه لغرض
البرهان بالتناقض سالباً. لكنك تستطيع في كثير من الأحيان إعادة صياغة قضية
موجبة بسهولة على هيئة قضية سالبة، وقضية سالبة على هيئة قضية موجبة.
ولكان برهاننا السابق هو نفسه في جوهره لو أثبتنا «$A = \emptyset$» بدلاً
من التالي السالب «لا توجد في $A$ !!{element}». (بحسب تعريف $=$، تعد
«$A = \emptyset$» قضية كلية، إذ تُفك إلى «كل !!{element} في~$A$ هي
!!{element} في~$\emptyset$، والعكس».) ويسهل أن نرى أنها تكافئ القضية
السالبة «ليس صحيحاً أنه يوجد $x \in A$».

لذلك يسهل أحياناً استعمال $\lnot p$ فرضاً أكثر من إثبات~$p$ مباشرة.
وحتى عندما يكون البرهان المباشر بالسهولة نفسها أو أسهل، كما في الأمثلة
التالية، يفضل بعض الناس السير بطريق غير مباشر. وإذا أربكك النفي المزدوج،
ففكر في البرهان بالتناقض على قضية ما بوصفه اشتقاق تناقض من القضية
\emph{المقابلة}. فبرهان $\lnot p$ بالتناقض هو اشتقاق تناقض من الفرض
$p$؛ وبرهان~$p$ بالتناقض هو اشتقاق تناقض من~$\lnot p$.

\begin{prop}
$A \subseteq A \cup B$.
\end{prop}

\begin{proof}
  نريد أن نبين أن $A \subseteq A \cup B$.
  \begin{quote}
    تبدو هذه، ظاهرياً، قضية موجبة: كل $x \in A$ ينتمي أيضاً إلى
    $A \cup B$. ونفيها هو: يوجد $x \in A$ لا ينتمي إلى $A \cup B$.
    فنستطيع إثبات القضية بطريق غير مباشر، بأن نفترض هذه القضية المنفية
    ونبين أنها تؤدي إلى تناقض.
    \end{quote}
  افترض العكس، أي $A \nsubseteq A \cup B$.
  \begin{quote}
    لدينا تعريف $A \subseteq A \cup B$: كل $x \in A$ ينتمي أيضاً إلى
    $A \cup B$. ولكي نفهم معنى $A \nsubseteq A \cup B$، علينا إجراء
    بعض المعالجة المنطقية الأولية للتعريف المفكوك: كون انتماء كل
    $x \in A$ إلى $A \cup B$ كاذباً يكافئ وجود \emph{بعض}
    $x \in A$ بحيث $x \notin A \cup B$. (وهذا موضع ينبغي فيه توخي
    الحذر الشديد؛ فكثير من محاولات الطلاب للبرهان بالتناقض تفشل بسبب
    تحليلهم الخاطئ لنفي قضية من قبيل «كل $A$ هو $B$».) وبعبارة أخرى،
    $A \nsubseteq A \cup B$ إذا وفقط إذا وجد $x$ بحيث $x \in A$ و
    $x \notin A \cup B$. وما بعد ذلك يسير.
  \end{quote}
  إذن يوجد $x \in A$ بحيث $x \notin A \cup B$. وبحسب تعريف $\cup$،
  يكون $x \in A \cup B$ إذا وفقط إذا كان $x \in A$ أو $x \in B$.
  وبما أن $x \in A$، فإن $x \in A \cup B$. وهذا يناقض الفرض
  $x \notin A \cup B$.
\end{proof}

\begin{prob}
أثبت \emph{بطريق غير مباشر} أن $A \cap B \subseteq A$.
\end{prob}

\begin{prop}
إذا كان $A \subseteq B$ و$B \subseteq C$ فإن $A \subseteq C$.
\end{prop}

\begin{proof}
  افترض أن $A \subseteq B$ و$B \subseteq C$. نريد أن نبين أن
  $A \subseteq C$.
  \begin{quote}
    لنسلك طريقاً غير مباشر: نفترض نفي ما نريد إثباته.
  \end{quote}
  افترض العكس، أي $A \nsubseteq C$.
  \begin{quote}
    كما سبق، نستدل على أن $A \nsubseteq C$ إذا وفقط إذا لم يكن كل
    $x \in A$ منتمياً أيضاً إلى $C$، أي إذا كان بعض $x \in A$ لا ينتمي
    إلى $C$. ولا تقلق؛ فمع الممارسة لن تحتاج بعد الآن إلى كثير من التفكير
    لفك نفي كهذا.
  \end{quote}
  وبعبارة أخرى، يوجد~$x$ بحيث $x \in A$ و$x \notin C$.
  \begin{quote}
    نستطيع الآن استعمال ذلك للوصول إلى التناقض. وبالطبع سنحتاج إلى
    الفرضين الآخرين.
  \end{quote}
  بما أن $A \subseteq B$، فإن $x \in B$. وبما أن $B \subseteq C$،
  فإن $x \in C$. لكن هذا يناقض $x \notin C$.
\end{proof}

\begin{prop}
إذا كان $A \cup B = A \cap B$ فإن $A = B$.
\end{prop}

\begin{proof}
  افترض أن $A \cup B = A \cap B$. نريد أن نبين أن $A = B$.
  \begin{quote}
    أصبحت البداية الآن مألوفة:
  \end{quote}
  افترض، للوصول إلى تناقض، أن $A \neq B$.
  \begin{quote}
    فرضنا في البرهان بالتناقض هو $A \neq B$. ولما كان $A = B$ إذا
    وفقط إذا كان $A \subseteq B$ و$B \subseteq A$، فإن
    $A \neq B$ إذا وفقط إذا كان $A \nsubseteq B$ \emph{أو}
    $B \nsubseteq A$. (لاحظ مقدار أهمية الحذر عند معالجة النفي!{})
    ولاشتقاق تناقض من هذا الفصل نستعمل برهاناً بالحالات ونبين أن تناقضاً
    ينتج في كل حالة.
  \end{quote}
  يكون $A \neq B$ إذا وفقط إذا كان $A \nsubseteq B$ أو
  $B \nsubseteq A$. ونميز حالتين.
  \begin{quote}
    في الحالة الأولى نفترض $A \nsubseteq B$، أي إنه يوجد $x$ بحيث
    $x \in A$ ولكن $x \notin B$. وتُعرَّف $A \cap B$ بأنها مجموعة
    !!{element} المشتركة بين $A$ و$B$؛ فإذا لم ينتم شيء إلى إحداهما لم
    ينتم إلى تقاطعهما. أما $A \cup B$ فهو $A$ مع $B$، ولذلك ينتمي إلى
    الاتحاد كل ما ينتمي إلى إحداهما. وهذا يخبرنا بأن
    $x \in A \cup B$ ولكن $x \notin A \cap B$، ومن ثم
    $A \cap B \neq A \cup B$.
  \end{quote}

  الحالة 1: $A \nsubseteq B$. عندئذ يوجد $x$ بحيث $x \in A$ ولكن
  $x \notin B$. وبما أن $x \notin B$، فإن $x \notin A \cap B$.
  وبما أن $x \in A$، فإن $x \in A \cup B$. إذن
  $A \cap B \neq A \cup B$، خلافاً للفرض
  $A \cap B = A \cup B$.

  الحالة 2: $B \nsubseteq A$. عندئذ يوجد $y$ بحيث $y \in B$ ولكن
  $y \notin A$. وكما سبق، $y \in A \cup B$ ولكن
  $y \notin A \cap B$، ومن ثم $A \cap B \neq A \cup B$، وهذا يناقض
  مرة أخرى $A \cap B = A \cup B$.
\end{proof}

\end{document}
